\chapter{Functions}
\label{ch:functions}

\section{Intro}

Functions are code blocks that can be invoked at a later time.

In Giovi, functions are first-class citizens, meaning that they are like any
other value: they have a well defined type (signature), they can be assigned to
variables, returned from other functions or passed to them as arguments, and so
on.

\paragraph{Function signature} A function signature is made of the function
name, its formal parameters (if any) and its return type (if any).

\paragraph{Declaration} A "traditional" function can be declared
by specifying its signature and its body.

\begin{minted}{\giovilexer}
buildGreetingMessage(name Text) => Text {
	if name = "simonvic" do {
		=> "heyo $name, how are you?"
	}
	=> "greetings, $name."
}
\end{minted}

It can later be invoked.

\nb{You have to declare the function first before invoking it. Hoisting is
still an experimental feature.}

\begin{minted}{\giovilexer}
buildGreetingMessage("simonvic") -> _out
\end{minted}

A function can also exist, for example, as a variable, where its type is the
signature. In this case the variable is an "empty function"; before invoking
it, you want to assign an actual function instance (with compatible signature.

\begin{minted}{\giovilexer}
greetingBuilder function(name Text) => Text
greetingBuilder <- buildGreetingMessage
# greetingBuilder is now "referencing" the previously
# defined buildGreetingMessage function
greetingBuilder("simonvic") -> _out
\end{minted}

\section{The dotref operator}

A "traditional" function declaration is semantically similar to a declaration
of a constant of type function, but with a difference: while there cannot be
two equally named functions defined in the same scope, there can be a function
and a variable (or constant) sharing the same name.

However, a function is always shadowed by a variable (or constant) with the
same name. To explicitly refer to the function, the \code{::} prefix is
required.

\begin{minted}{\giovilexer}
var foo <- 69
function foo() { "hello world" -> _out }
foo -> _out # 69
::foo -> _out # function(){}
\end{minted}

The dotref operator can also be used instead of a dot operator

\begin{minted}{\giovilexer}
struct SomeStruct <
	foo Int default 69
	function foo() { "hello world" -> _out }
>

someStruct SomeStruct
someStruct.foo -> _out # 69
someStruct::foo -> _out # function(){}
\end{minted}

\clearpage

\section{Function scope}

The scope of a function behaves like a normal scope, which means that, unlike
many other languages, it doesn't have a "call stack". It can be seen as a
normal scope in which some variables (i.e. formal parameters) are pre-declared
and declarations in parent scopes are visible.

\begin{minted}{\giovilexer}
const DEG_TO_RAD <- 0.01745329
degreeToRadians(deg Int) => Int {
	=> deg * DEG_TO_RAD
}
degreeToRadians(180) -> _out # 3.14
\end{minted}

\nb{This behavior, while it offers a lot of flexibility, it has some caveats:
when a function, that depends on a declaration of a parent scope, is invoked
within a scope where such declaration is not present, the program will error.}

\begin{minted}{\giovilexer}
degreeToRadians(deg Int) => Int {
	=> deg * DEG_TO_RAD
}
if true do {
	const DEG_TO_RAD <- 0.01745329
	degreeToRadians(180) -> _out # 3.14
} else {
	degreeToRadians(180) -> _out # ERROR: DEG_TO_RAD not declared
}
\end{minted}

\subsection{Nested functions}

A function declaration can appear anywhere, even within other function scopes.
Like any other declaration, its lifetime is bound to its enclosing scope,
meaning that after its enclosing scope has ended, the function will not be
visible anymore.

This can be very useful when you have repeating code, but you don't want to
expose a function.

\begin{minted}{\giovilexer}
function makePizza(ingredients [Text]) => Pizza {
	if "pineapple" in ingredients do panic()

	function cook(ingredient Text, seconds Int) {
		# cook the ingredient ...
	}

	cookedIngredients [Text]
	for ingredient in ingredients do {
		if ingredient {
			= "tomato" do cook(ingredient, 60) ->+ cookedIngredients
			= "fries" do cook(ingredient, 30) ->+ cookedIngredients
		} else {
			cook(ingredient, 100) ->+ cookedIngredients
		}
	}
	=> assemblePizza(cookedIngredients)
}
\end{minted}

\clearpage

\section{Parameters and invocation}

A formal parameter is the parameter that is declared in the function signature;
it has a name and a type. A function can have zero or more formal parameters.
They share the variable declaration syntax since they could be seen as such;
when a function is invoked, before executing the first statement of its body, a
new function scope is created and the formal parameters are declared as
variables.

The actual parameters (a.k.a. arguments) of the invocation are passed by value,
which means they're copied, therefore, while the function body can mutate the
parameters, the mutation will not be visible outside the function.

\begin{minted}{\giovilexer}
someFunction(someParameter Int) {
	someParameter -> _out # 69
	someParameter <- 420
	someParameter -> _out # 420
}
someArgument Int <- 69
someFunction(someArgument)
someArgument -> _out # still 69
\end{minted}

\todo{
	menzionare eccezione per self functions, se decido di "formalizzare il bug" :)

	Alternativamente, formalizzare che le strutture (o tutti i prodotti tipi)
	sono mutabili
}

\paragraph{Parameters defaults}

A formal parameter, being like a normal variable, also has a default value. You
can pass a default value by simply using the \code{default} keyword or the
underscore \code{\_}

\begin{minted}{\giovilexer}
someFunction(x Int, y Text) {
	x -> _out # 0
	y -> _out # "" (empty string)
}
someFunction(default, default)
someFunction(_, _)
\end{minted}

\subsection{Optional parameters}

Giovi supports optional parameters: by overriding a formal parameter's default
value, the parameter will be marked as optional. An optional parameter can be
omitted at invocation; its default value will be used instead.

\begin{minted}{\giovilexer}
someFunction(x Int default 69, y Text default "hello world") {
	x -> _out # 69
	y -> _out # hello world
}
someFunction()
\end{minted}

You can mark a parameter as optional without overriding its default

\begin{minted}{\giovilexer}
someFunction(x Int default _) { }
# equivalent to
someFunction(x Int default Int.default) { }
\end{minted}

It's possible to mix optional and required parameters, however only trailing
(last) optional parameters can be omitted when invoking using the positional
approach (\see{sec:arg_approach})

\begin{minted}{\giovilexer}
processPizzaOrder(
  kind Text default "Margherita",
  tableNumber Int,
  withPineapple Bool default false,
) {
	if withPineapple do panic()
	makePizza(kind) -> pizza Pizza
	serve(pizza, tableNumber)
}

processPizzaOrder("Diavola", 69) # "Diavola" to table 69
processPizzaOrder(_, 69) # "Margherita" to table 69
processPizzaOrder(_, _) # "Margherita" to table 0
processPizzaOrder(_, 69, true) # "Margherita" with pineapple to table 69
processPizzaOrder(69) # ERROR: expecting pizza type, but got 69
\end{minted}

\subsection{Positional and nominal parameters}
\label{sec:arg_approach}

\paragraph{Positional parameters} When invoking a function with the positional
approach, you just pass the values as arguments, which will be assigned to the
respective formal parameters in the same order: the first argument is assigned
to the first formal parameter, the second to the second and so on.

\paragraph{Nominal parameters} \label{sec:nominal_params} Giovi supports
nominal arguments: you can specify the name of the argument. When using the
nominal approach, the order of the arguments is not important.

\begin{minted}{\giovilexer}
processPizzaOrder(kind Text, tableNumber Int, withPineapple Bool) { #{...}# }
processPizzaOrder(tableNumber:69, withPineapple:false, kind:"Diavola")
\end{minted}

When using the nominal approach, all optional parameters can be omitted, as
long as the required parameters are specified.

\begin{minted}{\giovilexer}
processPizzaOrder(
  kind Text default "Margherita",
  tableNumber Int,
  withPineapple Bool default false
) { #{...}# }
processPizzaOrder(tableNumber:69) # Margherita to table 69 and no pineapple
\end{minted}

\subsection{Invocation labels}

As we've seen in \See{sec:nominal_params}, parameters names can be specified at
invocation.

When declaring a function, you sometimes have to decide whether to name a
parameter to either favor readability in the function body at the expense of
invocation or viceversa.

% \begin{minted}{\giovilexer}
% processPizzaOrder(
%   kind Text,
%   tableNumber Int,
%   hasPineapple Bool,
% ) {
% 	if hasPineapple do return
% 	makePizza(kind) -> pizza Pizza
% 	serve(pizza, tableNumber)
% }
%
% processPizzaOrder(
%   kind: "Diavola",
%   tableNumber: 69,
%   hasPineapple: false,
% )
% \end{minted}
%
% ... or viceversa
%
% \begin{minted}{\giovilexer}
% processPizzaOrder(
%   ofKind Text
%   atTableNumber Int,
%   withPineapple Bool,
% ) {
% 	if withPineapple do return
% 	makePizza(ofKind) -> pizza Pizza
% 	serve(pizza, atTableNumber)
% }
%
% processPizzaOrder(
%   ofKind: "Diavola",
%   atTableNumber: 69,
%   withPineapple: false,
% )
% \end{minted}

To get the best of both worlds and further improve readability, Giovi allows
you to define an "invocation label". Such label is the name of the parameter
that is exposed to the invocation, while the formal parameter name remains
unchanged in the function body. If an invocation label is not specified, the
formal parameter name will be used.

\begin{minted}{\giovilexer}
processPizzaOrder(
  ofKind: kind Text,
  atTable: tableNumber Int,
  hasPineapple Bool,
) {
	if hasPineapple do panic()
	makePizza(kind) -> pizza Pizza
	serve(pizza, tableNumber)
}

processPizzaOrder(
  ofKind: "Diavola",
  atTable: 69,
  hasPineapple: false
)
\end{minted}

\clearpage

\section{Struct functions}
\label{sec:struct_functions}

Functions can be defined within structs (\see{sec:structs}) and will be
associated to the type.\\*
Beware that these are, by definition, NOT methods; it can be seen as simple
syntax sugar to better organize code using types as "namespaces" for these
functions. We could make an analogy with "static methods" in OO languages.

\begin{minted}{\giovilexer}
Vehicle <
	speed Int

	accelerate(v Vehicle, amount Int) {
		amount ->+ v.speed
	}

	brake(v Vehicle) {
		v.speed -<- 1
	}

	speedToMilesPerHour(v Vehicle) => Int {
		=> v.speed / 1.6093
	}
>

car Vehicle
Vehicle.accelerate(car, 10)
"${car.speed} KMH" -> _out
"${Vehicle.speedToMilesPerHour(car)} MPH" -> _out
\end{minted}

Note that, although functions can be invoked from an instance of a struct,
they're still only related to the type; the syntax is allowed because of self
functions (\see{sec:self_func}).

\subsection{Self functions}
\label{sec:self_func}

More often than not, functions defined within a struct type are meant to
operate on such type (e.g. creating, updating, reading).

The \code{self} keyword can be prefixed to the function declaration; it's a
shortcut for defining a function that has a first parameter of the same type of
the struct in which the function is defined.

\begin{minted}{\giovilexer}
Vehicle <
	speed Int
	self speedToMilesPerHour() => Int {
		=> self.speed / 1.6093
	}
>
\end{minted}

Albeit \code{self} is a reserved keyword and cannot be used as parameter name,
the previous example would be equivalent to the following

\begin{minted}{\giovilexer}
Vehicle <
	speed Int
	speedToMilesPerHour(self Vehicle) => Int {
		=> self.speed / 1.6093
	}
>
\end{minted}

This also allows you to invoke the function using the instance of the struct as
the left-hand side of the dot operator, which will be passed automatically as
the first argument.

\begin{minted}{\giovilexer}
car Vehicle
car.speedToMilesPerHour() -> _out
# equivalent to
Vehicle.speedToMilesPerHour(car) -> _out
\end{minted}

\subsection{Self returning functions}

The \code{self} keyword can also be used as a return type as a shorthand.\\*
This comes very handy for writing "chainable" functions.

\begin{minted}{\giovilexer}
Color <
	r Int
	g Int
	b Int

	rgb(r Int, g Int, b Int) => self {
		=> Color.<r g b>
	}

	self brigthen() => self {
		self.r *<- 0.5
		self.g *<- 0.5
		self.b *<- 0.5
		=> self
	}

	self saturate() => self {
		# ...
		=> self
	}
>

Color.rgb(5, 20, 512)
	.brigthen()
	.saturate()
	.brighten()
\end{minted}

